% Compile with pdflatex + bibtex (Overleaf: Menu > Compiler > pdfLaTeX).
\documentclass[acmsmall,nonacm,screen]{acmart}

\usepackage{booktabs}
\usepackage{graphicx}

\begin{document}

\title{Does Automated Code Formatting Reduce Formatting-Related Code Churn?}
\subtitle{An Empirical Study of Python Repositories}

\author{Andriy Frolov}
\affiliation{%
  \institution{Technical University of Ko\v{s}ice}
  \city{Ko\v{s}ice}
  \country{Slovakia}}

\begin{abstract}
TODO: write last. Context, gap, method (sample size, classifier), main findings with numbers, implication.
\end{abstract}

\keywords{code formatting, Black, code churn, mining software repositories, Python}

\maketitle

\section{Introduction}
\label{sec:intro}
% Motivation: formatting differences in diffs and reviews; formatters promise to remove them.
% Gap: no evidence on what happens to formatting churn after adoption.
% Research questions RQ0--RQ3 and contributions.

\section{Background and Related Work}
\label{sec:related}
% Black and other formatters; code churn; classifying refactoring/formatting changes;
% MSR studies on tool adoption; \cite{muzammil2026hotspots}.

\section{Study Design}
\label{sec:design}
\subsection{Repository Selection}
\subsection{Project Start and Black Adoption Dates}
\subsection{Classifying Changes}
\subsection{Measures and Analysis}

\section{Classifier Validation}
\label{sec:validation}

\section{Results}
\label{sec:results}
\subsection{RQ0: How Do Python Projects Adopt Black?}
\subsection{RQ1: Formatting-Related Churn Before and After Adoption}
\subsection{RQ2: Formatting-Only Commits}
\subsection{RQ3: Evolution After Adoption}

\section{Discussion}
\label{sec:discussion}

\section{Threats to Validity}
\label{sec:threats}
\subsection{Construct Validity}
\subsection{Internal Validity}
\subsection{External Validity}
\subsection{Reliability}

\section{Conclusion}
\label{sec:conclusion}

\section*{Data Availability}
Scripts, the study protocol, and derived data are available at
\url{https://github.com/ryuuseiaika888/black-formatting-churn}.

\bibliographystyle{ACM-Reference-Format}
\bibliography{references}

\end{document}
